\section{Live Tournament Results}
\label{sec:live-results}

This section reports how the system performed in live deployment
during the 2026 World Cup. We interleave two records. The first is a
decision log of our actual fantasy entry: for each round we state the
recommendation produced by the models, the decision finally made, the
realised score, and a short post-hoc diagnosis. The second is a
retrospective cross-model backtest that replays every backend
end-to-end on frozen pre-round data, separating the question of how
each backend performs as an autonomous system from how the deployed
entry, which mixed model recommendations with domain judgement,
performed. Table~\ref{tab:backtest-rounds} gives per-round realised
points for both records, Figures~\ref{fig:backtest-cumulative}
and~\ref{fig:backtest-rounds} plot the cumulative and per-round
trajectories, and Table~\ref{tab:season-summary} tracks the official
(ensemble) recommendation against the actual entry and a random-squad
baseline.

\subsection{Cross-model backtest protocol}
\label{sec:live-protocol}

For each completed round we load the feature snapshot closest in
time to, but strictly before, the round's transfer deadline, run
every backend on that snapshot, and pass each backend's predicted
points to the MILP optimiser of Section~\ref{sec:methodology}, which
selects a 15-player squad and an 11-player starting eleven under the
stage's budget and country-cap constraints. The selected squad is
then scored with realised round points from the latest collector
snapshot. As a floor reference we also draw 1{,}000 budget- and
constraint-valid random squads per round and report the mean of
their realised scores. The replay script ships with the repository
and runs with a fixed random seed.

One caveat governs every comparison in this section: the backtest
grants each backend unlimited free transfers every round, whereas
the live entry is limited to two free transfers per stage
transition, with a three-point penalty per extra transfer. Backend
totals are therefore an upper bound on standalone model performance
rather than a like-for-like comparison with the actual entry, whose
figures are net of all transfer penalties incurred. A
constrained-transfer replay is deferred to
Section~\ref{sec:future_work}.

\subsection{Group stage decision log}
\label{sec:live-group}

\paragraph{Matchday 1: initial squad.}
The pre-tournament squad was assembled from the backends'
recommendations combined with domain priors, with the heuristic
backend as the primary live signal and the Poisson and GBM backends
as cross-checks. The pre-tournament call to include Willian Pacho
(Ecuador, defender, \$4.4M) came from the second author's match
observation; Pacho went on to return 11 points across the group
stage at 5.8\% ownership, one of the squad's best returns per
million. We captained Lautaro Mart\'{\i}nez (Argentina, forward,
\$8.8M, 3.2\% ownership) in a 4-4-2. The round realised 33 points,
last place in our personal league. Three failure modes dominate the
diagnosis. First, the captain blanked: Mart\'{\i}nez returned 1 point,
doubled to 2, and the captaincy multiplier is the largest single
variance source in a round. Second, a France triple-stack (Dou\'e,
Olise, and Demb\'el\'e at \$7.5M, \$9.5M, and \$10.0M) shared output
rather than concentrating it, returning 12 points across three
premium midfield slots. Third, none of the four defenders kept a
clean sheet, and Argentina's clean sheet went unexploited because no
Argentine defender was paired with Emiliano Mart\'{\i}nez in goal.

\paragraph{Matchday 2.}
With two free transfers available, the optimiser's best two-transfer
move was to sell Olise and Pacho for Saka and Iba\~nez. We rejected
the Olise sale on domain grounds: Olise had been the squad's top
Matchday~1 scorer with 6 points and plays for a tournament
favourite. The deciding input was the second author's read on
Olise's form, and the captaincy was moved from Mart\'{\i}nez, who was
transferred out, to Olise. The round realised 84 points net, the
entry's largest group-stage total. Olise returned 9 raw points,
doubled to 18 as captain, the squad's largest captain return of the
group stage; Demb\'el\'e's 12 points justified retaining the
remaining France pair; Gakpo added 19; and James and Pacho
contributed 9 each in defence. (An earlier internal draft recorded
this round as 102 points; the realised total in the scoring artifact
is 84, and we report realised figures throughout.) The post-hoc
reading is symmetric: the optimiser was right that Olise carried
future ceiling, and the domain override was right that selling him
was wrong. For a premium player in form, anchoring on form beat the
model's expected value by roughly the margin the captaincy
multiplier implies.

\paragraph{Matchday 3.}
Both free transfers were used and two additional paid transfers were
taken at a cost of six points, bringing in Daniel Mu\~noz, Alexander
Freeman, Ronald Ara\'ujo, and Deniz Undav while keeping Messi as
captain in a 5-2-3. The central risk in the final group round is
rotation, since teams that have already clinched their group rest
starters. The rotation-risk flags for this round, including which
clinched teams would rotate, again came from the second author's
reading of observed squad management; the model carried only a
country-level multiplier and proved too conservative. Selling Pacho
ahead of Ecuador's fixture against Mexico was correct in spirit:
Mexico won and Pacho's defensive return was modest. The round
realised 46 points gross and 40 net of the transfer penalty (the
in-app display showed 42; we note the two-point reconciliation for
completeness and report the artifact figure). Captain Messi returned
7 raw points, doubled to 14, against Jordan: a defensible pick,
although in retrospect Demb\'el\'e outscored him within the same
starting eleven. The paid transfers returned roughly the six points
they cost, a net-zero outcome consistent with the unforced-hit cost
discussed in Section~\ref{sec:lessons}.

\subsection{Knockout rounds}
\label{sec:live-knockout}

\paragraph{Round of 32: wildcard rebuild.}
We played the wildcard, which permits unlimited transfers without
penalty, and rebuilt the squad almost from scratch from the
post-group candidate pool: Rangel and Beach in goal; Laporte,
Cucurella, Mu\~noz, Nuno Mendes, and Freeman in defence;
Demb\'el\'e, Vin\'{\i}cius, Bellingham, Olise, and Manzambi in
midfield; Messi, Haaland, and Undav up front. Total cost was
\$102.9M against the \$105.0M knockout-stage cap. The optimiser
marginally preferred Messi as captain; we chose Demb\'el\'e on
differential grounds, at 19.7\% ownership against Messi's 37.8\%.
The realised outcome of the round appears in
Table~\ref{tab:backtest-rounds}.

\paragraph{Round of 16 and quarter-final.}
The entry scored 66 points net in the Round of 16 and 73 in the
quarter-final, in-app records taken as primary (the squad sheets for
these two rounds were not archived in the repository, an operational
lapse we note for completeness; every other round's squad is
registered pre-deadline). Both rounds continued the wildcard squad's
France--Spain--England spine under the four-transfer knockout quotas.
The quarter-final was the entry's second-best round of the
tournament and beat every backend replay except the GBM's 76
(Table~\ref{tab:backtest-rounds}): the ensemble replay managed 64,
the random baseline 22.5.

\paragraph{Semi-final.}
The semi-final decision is the best-documented of the tournament
because it was registered as a pre-round head-to-head against the
model (the artifact behind Section~\ref{sec:analysis}'s override
analysis). Entering the round with two free transfers and \$1.0M in
the bank, the optimiser's remaining recommendation was to sell
Mac~Allister for Dou\'e and a dead bench slot for Spence; the
operator declined both and kept the squad. The captaincy went to
Mbapp\'e (vice Messi) with the Clean Sheet Shield booster active,
and the pre-round simulation put the operator's pick within one
expected point of the model's (59.7 versus 60.5). The round realised
37 points net, the worst of the knockouts: France lost the
semi-final 0--2 to Spain and the entry's France triple stack
returned four raw points between Maignan, Upamecano, Demb\'el\'e,
and the captain, whose doubled score was 2. The two Spain defenders
(Laporte 7, Cucurella 6) and Messi's 8 kept the round from single
digits, and the Shield converted Argentina's single goal conceded
into one extra point for Mac~Allister, reconciling the recorded 37
with the 36 an id-based rescore produces. The backend replays fared
little better (ensemble 37, GBM 59); with only two fixtures, a
semi-final round's realized total carries high variance under any
selection policy. After the semi-finals the entry stood
at \GlobalRankSF{} in the global game, on 418 points net across the
seven completed rounds.
\paragraph{Final round.}
Round 8 covered both closing matches (the third-place match and the
final), and the deadline was the bronze kickoff, so the round was
planned as a whole. The locked plan moved five transfers into the
eventual champion's block, where Sim\'on, Porro, Cubars\'i, and
Baena joined Laporte and Cucurella for six Spain starters, with the
Qualification Booster active and Mbapp\'e captaining the bronze
match. One further slot went to a second England midfielder as
bronze-game cover, a buy the model rated below its Spain
alternatives and that returned nothing. A planned mid-round armband switch to Messi had to be
abandoned when live play established that the app locks the
captaincy at the captain's own kickoff, contrary to the help-page
wording (the rules transcription and the live tooling have been
corrected accordingly); the operator therefore committed a static
5--3--2 before the bronze kickoff, benching Kane on a rotation read
and starting a fifth defender. Both calls at the operator's
discretion were correct: Kane never played, and the model's one
clear error was the omission of Olise (10 points, started and
scored) from the model-guided eleven. The forecaster's record on
the two matches was one correct call and one miss: the final
matched the modal forecast outcome (Spain 1--0, forecast 54\%),
while the bronze ended England 6--4 over
the 59\%-favoured France, the highest-scoring match of the
tournament and a forecast miss, though Mbapp\'e's hat-trick within
the loss returned 17 raw points, doubled to 34 as captain, the
entry's largest single return of the tournament. The round realised \FinalRoundNet{} points net,
the entry's highest of the eight rounds: the six Spain starters
returned the clean sheet plus the booster's twelve points, and the
season closed at \TotalUserActualNet{} net points and global rank
\GlobalRankFinal{}, an improvement of approximately 250{,}000
places from \GlobalRankSF{} after the semi-finals.

\subsection{Cross-model findings}
\label{sec:live-findings}

Across the seven completed rounds (Table~\ref{tab:backtest-rounds}
and Figure~\ref{fig:backtest-cumulative}), the GBM backend leads the
cumulative leaderboard with \TotalGbm{} points, followed by the
per-position ensemble with \TotalEnsemble{}. Heuristic~v2
(\TotalHeuristicVTwo{}), Monte Carlo (\TotalMonteCarlo{}), and the
original heuristic (\TotalHeuristic{}) form a middle tier, while the
Poisson backend trails the field at \TotalPoisson{}. Every backend
clears the random baseline, whose cumulative mean is
\TotalRandomBaselineMean{} points. Our actual entry stands at
\TotalUserActualNet{} points net of penalties, which is not directly
comparable to the unlimited-transfer backend totals but sits between
the middle tier and the leaders despite paying real transfer costs.
Three observations follow.

First, the Poisson deficit is a genuine mis-calibration finding
rather than noise: the backend systematically over-predicts mean
realised points, and its squads skew toward high-volume,
low-realisation players. We do not recommend deploying it in its
current form, despite its held-out goalkeeper RMSE win
(Section~\ref{sec:validation}), which we now suspect reflects a
small per-position sample; a recalibration pass is queued in
Section~\ref{sec:future_work}.

Second, the captaincy is the highest-leverage single decision in a
round, and it resists prospective optimisation. In each of the three
group rounds the captain that proved optimal within our own starting
eleven was identified only in retrospect: Gakpo over Olise in
Matchday~2 and Demb\'el\'e over Messi in Matchday~3. Per-round
captain outcomes are visible in the round-level spread of
Figure~\ref{fig:backtest-rounds}.

Third, the rounds that went well combined the model signal with
domain judgement (the Olise captaincy, the Pacho pick, the rotation
flags), while the rounds that went poorly weighted one signal over
the other. Table~\ref{tab:season-summary} tracks this round by
round against the official ensemble recommendation.

The closing leaderboard over all eight rounds reads: GBM
\TotalGbm{}, ensemble \TotalEnsemble{}, heuristic~v2
\TotalHeuristicVTwo{}, the actual entry \TotalUserActualNet{} net,
Monte Carlo \TotalMonteCarlo{}, heuristic \TotalHeuristic{}, Poisson
\TotalPoisson{}, random baseline \TotalRandomBaselineMean{}. The
headline comparison is unchanged from mid-tournament but now
complete: the strongest learned backend beats the deployed hybrid
entry under unlimited free transfers, the hybrid entry beats half
the backend field while paying real transfer costs, and everything
beats chance by a factor of three or more. The entry's strongest
rounds (Matchday~2's 84 and the final's \FinalRoundNet{}) were
precisely the rounds where model recommendations and domain
judgement pointed the same way.

\begin{table}[t]
  \centering
  \small
  % GENERATED by scripts/paper_tables.py -- do not edit
\begin{tabular}{lrrrrrrrr}
\toprule
Round & Heuristic & Heuristic v2 & Poisson & GBM & Monte Carlo & Ensemble & User (actual) & Random baseline \\
\midrule
Group Matchday 1 & 46 & 39 & 30 & 105 & 77 & 105 & 33 & 22 \\
Group Matchday 2 & 81 & 79 & 56 & 148 & 79 & 130 & 84 & 23 \\
Group Matchday 3 & 84 & 72 & 44 & 115 & 44 & 94 & 40 & 22 \\
Round of 32 & 73 & 87 & 42 & 108 & 81 & 104 & 83 & 22 \\
Round of 16 & 41 & 63 & 24 & 101 & 56 & 84 & 0 & 20 \\
\midrule
\textbf{Cumulative} & 325 & 340 & 196 & 577 & 337 & 517 & 240 & 105 \\
\bottomrule
\end{tabular}

  \caption{Realized fantasy points per round for every backend replayed
    on frozen pre-round data, alongside the user's actual entry and a
    random-squad baseline.
    \figsource{scripts/model\_backtest.py $\rightarrow$
    data/evaluation/backtest\_summary.json}}
  \label{tab:backtest-rounds}
\end{table}

\begin{figure}[t]
  \centering
  \includegraphics[width=0.85\textwidth]{fig_backtest_cumulative.pdf}
  \caption{Cumulative realized points per backend across the
    tournament.
    \figsource{scripts/model\_backtest.py}}
  \label{fig:backtest-cumulative}
\end{figure}

\begin{figure}[t]
  \centering
  \includegraphics[width=0.85\textwidth]{fig_backtest_rounds.pdf}
  \caption{Per-round realized points per backend.
    \figsource{scripts/model\_backtest.py}}
  \label{fig:backtest-rounds}
\end{figure}

\begin{table}[t]
  \centering
  % GENERATED by scripts/paper_tables.py -- do not edit
\begin{tabular}{lrrr}
\toprule
Round & Ensemble & User (actual) & Random baseline \\
\midrule
Group Matchday 1 & 105 & 33 & 22 \\
Group Matchday 2 & 130 & 84 & 23 \\
Group Matchday 3 & 94 & 40 & 22 \\
Round of 32 & 104 & 83 & 22 \\
Round of 16 & 84 & 0 & 20 \\
Quarter-final & \POSTFINAL{tbd} & \POSTFINAL{tbd} & \POSTFINAL{tbd} \\
Semi-final & \POSTFINAL{tbd} & \POSTFINAL{tbd} & \POSTFINAL{tbd} \\
Final & \POSTFINAL{tbd} & \POSTFINAL{tbd} & \POSTFINAL{tbd} \\
\bottomrule
\end{tabular}

  \caption{Season summary: official (ensemble) recommendation, user
    actual, and random baseline per round.
    \figsource{scripts/model\_backtest.py $\rightarrow$
    data/evaluation/backtest\_summary.json}}
  \label{tab:season-summary}
\end{table}
